%% kapa: exam=1 round=22 year=2011 period=2 form=A subject=appraisal_law subject_ko=부동산관계법규 seq=27 number=67
%% source: Q-Net 감정평가사 기출문제 (cst00309) artlSeq=5036790 file=extracted/1008/5036790/2011년 제22회 감정평가사 1차시험 기출문제/2교시.hwp
%% conversion: kordoc 4.24.51 -> tools/gen1.py (auto); review=proofread-gpt-astra; flags=none
\begin{question}{67}
국유재산법령상 국유재산으로 기부채납을 받을 수 있는 경우는?
\begin{choices}
  \item 국가에 기부하려는 재산이 재산가액 대비 유지ㆍ보수 비용이 지나치게 많은 경우
  \item 기부자의 상속인에게 무상으로 사용허가하여 줄 것을 조건으로 하여 국가에 행정재산으로 기부하는 경우
  \item 기부하려는 재산이 국가에 이익이 없는 것으로 인정되는 경우
  \item 기부자의 사망 후 상속인에게 반환하여 줄 것을 조건으로 국가에 재산을 기부하는 경우
  \item 특정한 행정재산의 용도를 변경하여 줄 것을 조건으로 국가에 재산을 기부하는 경우
\end{choices}
\end{question}
\answer{2}
